\section{L'action : histoire}

%G. Moreau

\subsection{Principe philosophique}
Principe philosophique d'économie naturelle : Aristote (- 4 ème), Saint
Augustin d'Hippone (4 ème) et Thomas d'Aquin (13 ème)
\subsection{17 ème Optique géométrique}
Héron d'Alexandrie (-1 er) : égalité des angles de réflexion.

Lois de W. Snell & R.Descartes [conservation de la quantité de mouvement
scalaire]

Principe de P. de Fermat [G. Leibniz, différrentielles]

"À l'interface entre deux milieux, la lumière se propage d'un point à un un
autre sur un trajet tel que la durée de parcours soit extrémale."

\[
\delta(\Delta t)=0 \qquad \Delta t \propto \int^{x_B}_{x_A}dx
\frac{c}{dx/dt_{onde}} = \int dx \frac{c}{v_O} =
\int dx \frac{c}{\lambda\nu} = \int dx \frac{pc}{E}
\]
avec
\[
\left{
\begin{array}{ c c c }
 p= & \hbar k & = \frac{h}{\lambda} \\
 E= & \hbar \omega & = \frac{h}{T} = h\nu \\
\end{array}
\]
(De Broglie & Planck-Einstein)

Jean Bernoulli trouve [calcul infinitésimal, "variationnel"] la véritable
brachistochrome : "lien à la mécanique"

EXPLICITE DUALITÉ ONDE-CORPUSCULE

(atomiste Démocrite d'Abdère (- 5 ème) corpularistes Isaac Newton,
calcul différentiel)

Description ondulatoire (lumière) - C. Huygens [conservation de la quantité
de mouvement sommée] Retrouve les lois de Snell-Descartes (+ A. Fresnel,
19 ème)

\subsection{}
\subsection{}
